% GENERATED FILE. Edit canonical lessons, then run npm run generate:books.
% canonical-source-hash: fnv1a64:d5683266128fbcc5
% canonical-lessons: ZH-C534-ningke, ZH-C534-jinliang, ZH-C534-jiaru, ZH-C534-xunuo, ZH-C534-wanshiruyi

\chapter{I'd Rather}
\label{ch:zh-i-d-rather}
\hlchaptermodality{534}

\begin{chapteropening}
\textbf{By the end of this chapter:} \emph{I can say I'd rather, that I'll do what I can, suppose something, give my word and wish someone all the best.}

Complete the last lesson of chapter 534: I can say I'd rather, that I'll do what I can, suppose something, give my word and wish someone all the best.
\end{chapteropening}

\section[nìngkě]{nìngkě — would rather}

\label{lesson:ZH-C534-ningke}



\begin{quote}
Before the new one: say the Mandarin for to give up all hope, then the Mandarin for may your wishes all come true.
\end{quote}

\subsection*{You'll want to know: nìngkě}
\emph{nìngkě} — \textquotedblleft{}would rather\textquotedblright{}.

\begin{quote}
\emph{wǒ nìngkě zǒulù yě bù xiǎng děng chē} — I'd rather walk than wait for the bus.
\end{quote}

Its characters come later; for now it is a word for the ear.

\begin{cousinweb}
\emph{nìngkě} is built from \emph{nìng} and \zh{可} 'may', both inherited from Middle Chinese, \emph{neng} and \emph{khaX} (Guangyun, in Baxter's transcription). The dictionary glosses the first 'calm' and, after an older dictionary, as a word of wishing. The falling tone of nìng is not the regular outcome of \emph{neng}, which gives níng.
\end{cousinweb}

\begin{grammarlens}[title={What you've built}]
One more word for saying what would happen.
\end{grammarlens}

\subsection*{Your turn}
\begin{itemize}
\raggedright
  \item \emph{Say it:} \emph{nìngkě}
  \item \emph{Say it:} \emph{nìngkě}, once more
  \item \emph{Say it:} \emph{wǒ nìngkě zǒulù yě bù xiǎng děng chē} — I'd rather walk than wait for the bus.
  \item \emph{Recall:} say the Mandarin for to give up all hope, then the Mandarin for may your wishes all come true, then say \emph{nìngkě} again
\end{itemize}

\begin{tcolorbox}[breakable,colback=teal!4,colframe=teal!35!black,title={Before you move on}]
What does nìngkě mean? (\textquotedblleft{}Would rather\textquotedblright{}.) Say it once more.
\end{tcolorbox}
\section[jǐnliàng]{jǐnliàng — as far as one can}

\label{lesson:ZH-C534-jinliang}



\begin{quote}
Before the new one: say the Mandarin for may your wishes all come true, then the Mandarin for would rather.
\end{quote}

\subsection*{You'll want to know: jǐnliàng}
\emph{jǐnliàng} — \textquotedblleft{}as far as one can\textquotedblright{}.

\begin{quote}
\emph{rúguǒ nǐ yǒu wèntí wǒ huì jǐnliàng bāng nǐ} — If you have a problem, I'll do what I can to help.
\end{quote}

Its characters come later; for now it is a word for the ear.

\begin{cousinweb}
\emph{jǐnliàng} is built from \emph{jǐn} 'to the full' and \zh{量} 'an amount'. No dictionary we use records an older form of \emph{jǐn} as this word writes it, so its history is not settled. \zh{量} is inherited from Middle Chinese \emph{ljangH} (Guangyun, in Baxter's transcription).
\end{cousinweb}

\begin{grammarlens}[title={What you've built}]
One more word for saying what would happen.
\end{grammarlens}

\subsection*{Your turn}
\begin{itemize}
\raggedright
  \item \emph{Say it:} \emph{jǐnliàng}
  \item \emph{Say it:} \emph{jǐnliàng}, once more
  \item \emph{Say it:} \emph{rúguǒ nǐ yǒu wèntí wǒ huì jǐnliàng bāng nǐ} — If you have a problem, I'll do what I can to help.
  \item \emph{Recall:} say the Mandarin for may your wishes all come true, then the Mandarin for would rather, then say \emph{jǐnliàng} again
\end{itemize}

\begin{tcolorbox}[breakable,colback=teal!4,colframe=teal!35!black,title={Before you move on}]
What does jǐnliàng mean? (\textquotedblleft{}As far as one can\textquotedblright{}.) Say it once more.
\end{tcolorbox}
\section[\texorpdfstring{\zh{假如} (jiǎrú)}{假如 (jiǎrú)}]{\zh{假如} (jiǎrú) — supposing that}

\label{lesson:ZH-C534-jiaru}



\begin{quote}
Before the new one: say the Mandarin for would rather, then the Mandarin for as far as one can.
\end{quote}

\subsection*{You'll want to know: \zh{假如}}
\textbf{\zh{假如}} — \emph{jiǎrú} — \textquotedblleft{}supposing that\textquotedblright{}.

\begin{quote}
\textbf{\zh{假如明天下雨我们就不去了}} \emph{jiǎrú míngtiān xiàyǔ wǒmen jiù bú qù le} — Supposing it rains tomorrow, we won't go.
\end{quote}

\begin{cousinweb}
\textbf{\zh{假如}} (\emph{jiǎrú}) is built from \zh{假} jiǎ and \zh{如} 'like', both inherited from Middle Chinese, \emph{kaeX} and \emph{nyo} (Guangyun, in Baxter's transcription), which glosses the first 'to borrow', 'false' and 'for the time being': a borrowed case.
\end{cousinweb}

\begin{grammarlens}[title={What you've built}]
One more word for saying what would happen.
\end{grammarlens}

\subsection*{Your turn}
\begin{itemize}
\raggedright
  \item \emph{Say it:} \emph{jiǎrú}
  \item \emph{Say it:} \emph{jiǎrú}, once more
  \item \emph{Say it:} \emph{jiǎrú míngtiān xiàyǔ wǒmen jiù bú qù le} — Supposing it rains tomorrow, we won't go.
  \item \emph{Recall:} say the Mandarin for would rather, then the Mandarin for as far as one can, then say \emph{jiǎrú} again
\end{itemize}

\begin{tcolorbox}[breakable,colback=teal!4,colframe=teal!35!black,title={Before you move on}]
What does \zh{假如} mean? (\textquotedblleft{}Supposing that\textquotedblright{}.) Say it once more.
\end{tcolorbox}
\section[xǔnuò]{xǔnuò — to give one's word; to pledge}

\label{lesson:ZH-C534-xunuo}



\begin{quote}
Before the new one: say the Mandarin for as far as one can, then the Mandarin for supposing that.
\end{quote}

\subsection*{You'll want to know: xǔnuò}
\emph{xǔnuò} — \textquotedblleft{}to give one's word; to pledge\textquotedblright{}.

\begin{quote}
\emph{rúguǒ nǐ xǔnuò le jiù yídìng yào zuò dào} — If you've given your word, you must keep it.
\end{quote}

Its characters come later; for now it is a word for the ear.

\begin{cousinweb}
\emph{xǔnuò} is built from \emph{xǔ} 'to allow' and \emph{nuò} 'to say yes', both inherited from Middle Chinese, \emph{xjoX} and \emph{nak} (Guangyun, in Baxter's transcription), which glosses the first 'to permit' and the second, after an older dictionary, 'to answer'.
\end{cousinweb}

\begin{grammarlens}[title={What you've built}]
One more word for saying what would happen.
\end{grammarlens}

\subsection*{Your turn}
\begin{itemize}
\raggedright
  \item \emph{Say it:} \emph{xǔnuò}
  \item \emph{Say it:} \emph{xǔnuò}, once more
  \item \emph{Say it:} \emph{rúguǒ nǐ xǔnuò le jiù yídìng yào zuò dào} — If you've given your word, you must keep it.
  \item \emph{Recall:} say the Mandarin for as far as one can, then the Mandarin for supposing that, then say \emph{xǔnuò} again
\end{itemize}

\begin{tcolorbox}[breakable,colback=teal!4,colframe=teal!35!black,title={Before you move on}]
What does xǔnuò mean? (\textquotedblleft{}To give one's word; to pledge\textquotedblright{}.) Say it once more.
\end{tcolorbox}
\section[\texorpdfstring{\zh{万事如意} (wànshì rúyì)}{万事如意 (wànshì rúyì)}]{\zh{万事如意} (wànshì rúyì) — may all go as you wish}

\label{lesson:ZH-C534-wanshiruyi}



\begin{quote}
Before the new one: say the Mandarin for supposing that, then the Mandarin for to give one's word.
\end{quote}

\subsection*{You'll want to know: \zh{万事如意}}
\textbf{\zh{万事如意}} — \emph{wànshì rúyì} — \textquotedblleft{}may all go as you wish\textquotedblright{}.

\begin{quote}
\textbf{\zh{希望你万事如意}} \emph{xīwàng nǐ wànshì rúyì} — I hope everything goes well for you.
\end{quote}

\begin{cousinweb}
\textbf{\zh{万事如意}} (\emph{wànshì rúyì}) is built from \zh{万} 'ten thousand', \zh{事} 'a matter', \zh{如} 'like' and \zh{意} 'what one means', all inherited from Middle Chinese, \emph{mjonH}, \emph{dzriH}, \emph{nyo} and \emph{'iH} (Guangyun, in Baxter's transcription): every matter as you would have it.
\end{cousinweb}

\begin{grammarlens}[title={What you've built}]
One more word for saying what would happen.
\end{grammarlens}

\subsection*{Your turn}
\begin{itemize}
\raggedright
  \item \emph{Say it:} \emph{wànshì rúyì}
  \item \emph{Say it:} \emph{wànshì rúyì}, once more
  \item \emph{Say it:} \emph{xīwàng nǐ wànshì rúyì} — I hope everything goes well for you.
  \item \emph{Recall:} say the Mandarin for supposing that, then the Mandarin for to give one's word, then say \emph{wànshì rúyì} again
\end{itemize}

\begin{tcolorbox}[breakable,colback=teal!4,colframe=teal!35!black,title={Before you move on}]
What does \zh{万事如意} mean? (\textquotedblleft{}May all go as you wish\textquotedblright{}.) Say it once more.
\end{tcolorbox}
